\subsection{整闭整环的分解与惯性}

引理 4. 设 A 为整闭整环，K 为其分式域，p 为 K 的特征指数，K' 为 K 的一个纯不可分扩张，A' 为 K' 中包含 A 且在 A 上整的子环。对 A 的每个素理想 p，存在唯一的素理想 \(\mathfrak{p}'\)，即 \(\mathbf{A}'\) 中位于 \(\mathfrak{p}\) 之上的素理想，且 \(\mathfrak{p}'\) 是由满足如下条件的 \(x \in \mathbf{A}'\) 组成的集合：存在整数 \(m \geqslant 0\) 使得 \(x^{p^m} \in \mathfrak{p}\)。

\(\mathfrak{p}'\) 的存在性由第 1 节定理 1 得出。若 \(x \in \mathfrak{p}'\)，则存在 \(m \geqslant 0\) 使 \(x^{p^m} \in \mathbf{K}\)，于是 \(x^{p^m} \in \mathbf{A}\)（因 \(\mathbf{A}\) 整闭），从而 \(x^{p^m} \in \mathfrak{p}' \cap \mathbf{A} = \mathfrak{p}\)。反之，若 \(x \in \mathbf{A}'\) 满足 \(x^{p^m} \in \mathfrak{p} \subset \mathfrak{p}'\)，则 \(x \in \mathfrak{p}\)（因 \(\mathfrak{p}'\) 素）。

注 (1). 由 § 1 第 3 节命题 11 的推论可知，A 在 \(\mathbf{K}'\) 中的整闭包是由满足如下条件的 \(x \in \mathbf{K}'\) 组成的集合：存在 \(m \geqslant 0\) 使 \(x^{p^m} \in \mathbf{A}\)（《代数》第 V 章 § 8 第 1 节命题 1）。

命题 6. 设 A 为整闭整环，K 为其分式域，K' 为 K 的拟 Galois 扩张，A' 为 A 在 K' 中的整闭包。那么：

\begin{enumerate}
\def\labelenumi{(\roman{enumi})}
\item
  对 A 的每个素理想 p，K' 的 K-自同构群 G 传递地作用于 A' 的位于 p 之上的素理想所成的集合。
\item
  对 \(\mathfrak{p}'\)（\(\mathbf{A}'\) 的任意素理想），分式域 \(k'\)（即 \(\mathbf{A}' / \mathfrak{p}'\) 的分式域）是 \(k\)（\(\mathbf{A} / (\mathbf{A} \cap \mathfrak{p}')\) 的分式域）的拟 Galois 扩张，且典范同态 \(\sigma \mapsto \bar{\sigma}\)（从 \(\mathcal{G}^{\mathbb{Z}}(\mathfrak{p}')\) 到群 \(\Gamma\)，后者由 \(k\)-自同构（\(k'\) 上的）组成）经取商给出 \(\mathcal{G}^{\mathbb{Z}}(\mathfrak{p}') / \mathcal{G}^{\mathbb{T}}(\mathfrak{p}')\) 到 \(\Gamma\) 上的一个双射。
\end{enumerate}

\begin{enumerate}
\def\labelenumi{(\Alph{enumi})}
\item
  先设 \(\mathbf{K}'\) 是 \(\mathbf{K}\) 的有限 Galois 扩张。此时 \(\mathbf{A} = \mathbf{A}' \cap \mathbf{K}\)（因 \(\mathbf{A}\) 整闭），故 \(\mathbf{A}\) 是 \(\mathcal{G}\) 在 \(\mathbf{A}'\) 中的不变量环。由于 \(\mathcal{G}\) 有限，命题在这一情形由第 2 节定理 2 得出。
\item
  其次设 \(\mathbf{K}'\) 是 \(\mathbf{K}\) 的任意 Galois 扩张。此时 \(\mathbf{K}'\) 是一个右定向族 \((\mathrm{K}_{\alpha})_{\alpha \in I}\)（由 \(\mathbf{K}\) 的有限 Galois 扩张组成）的并。为证 (i)，考虑 \(\mathfrak{p}'\)、\(\mathfrak{q}'\)，即 \(\mathbf{A}'\) 的位于 \(\mathfrak{p}\) 之上的两个素理想。对每个 \(\alpha \in I\)，\(\mathfrak{p}' \cap \mathrm{K}_{\alpha}\) 与 \(\mathfrak{q}' \cap \mathrm{K}_{\alpha}\) 是 \(\mathbf{A}' \cap \mathrm{K}_{\alpha}\) 的位于 \(\mathfrak{p}\) 之上的两个素理想。由于 \(\mathbf{A}' \cap \mathrm{K}_{\alpha}\) 是 \(A\) 在 \(\mathrm{K}_{\alpha}\) 中的整闭包，且在 \(\mathrm{K}_{\alpha}\) 上 \(\mathcal{G}\) 的元素的限制构成 \(\mathbf{K}\) 的 K-自同构群，故由情形 (A) 存在 \(\sigma \in \mathcal{G}\) 使 \(\sigma. (\mathfrak{p}' \cap \mathrm{K}_{\alpha}) = \mathfrak{q}' \cap \mathrm{K}_{\alpha}\)。设 \(\mathcal{E}_{\alpha}\) 为具有上述性质的 \(\sigma \in \mathcal{G}\) 组成的集。取 \(\sigma \in \mathcal{G} - \mathcal{E}_{\alpha}\)；则对每个 \(\tau \in \mathcal{G}\)，它使 \(\mathrm{K}_{\alpha}, (\sigma \tau). (\mathfrak{p}' \cap \mathrm{K}_{\alpha}) = \sigma. (\mathfrak{p}' \cap \mathrm{K}_{\alpha}) \neq \mathfrak{q}' \cap \mathrm{K}_{\alpha}\) 成立，因而 \(\sigma \tau \in \mathcal{G} - \mathcal{E}_{\alpha}\)。由此可知 \(\mathcal{E}_{\alpha}\) 在拓扑 Galois 群 \(\mathcal{G}\) 中是闭的（《代数》第 V 章附录 II 第 1 节），且族 \((\mathcal{E}_{\alpha})_{\alpha \in I}\) 显然是左定向的。由于 \(\mathcal{G}\) 紧（同前引，第 2 节，命题 3）且诸 \(\mathcal{E}_{\alpha}\) 非空，交 \(\mathcal{E}\)（即族 \((\mathcal{E}_{\alpha})\) 的交）非空，且有 \(\sigma. \mathfrak{p}' = \mathfrak{q}'\)（对每个 \(\sigma \in \mathcal{E}\)），由此得 (i)。
\end{enumerate}

为证 (ii)，注意 \(k'\) 是右定向族 \((k_{\alpha})_{\alpha \in \mathbf{I}}\) 的并，其中 \(k_{\alpha}\) 是 \((\mathrm{A}' \cap \mathrm{K}_{\alpha}) / (\mathfrak{p}' \cap \mathrm{K}_{\alpha})\) 的分式域。由于每个 \(k_{\alpha}\) 由 (A) 是 \(k\) 的拟 Galois 扩张，故 \(k'\) 亦然（《代数》第 V 章 § 6 第 3 节命题 8）。另一方面，设 \(u\) 是一个 \(k\)-自同构，作用于 \(k'\)，并设 \(\pi': \mathrm{A}' \to \mathrm{A}' / \mathfrak{p}'\) 为典范同态。对 \(\mathrm{A}' \cap \mathrm{K}_{\alpha}\) 应用第 2 节定理 2 可知，对每个 \(\alpha\) 存在一个非空集 \(\mathcal{F}_{\alpha}\)，由具有下述性质的元素 \(\sigma \in \mathcal{G}\) 组成：\(\sigma \cdot (\mathfrak{p}' \cap \mathrm{K}_{\alpha}) = \mathfrak{p}' \cap \mathrm{K}_{\alpha}\)，且 \(u(\pi'(x)) = \pi'(\sigma.x)\) 对一切 \(x \in \mathrm{A}' \cap \mathrm{K}_{\alpha}\) 成立。与上面一样可以看出 \(\mathcal{F}_{\alpha}\) 在 \(\mathcal{G}\) 中闭，且由于 \((\mathcal{F}_{\alpha})\) 是左定向集，其交 \(\mathcal{F}\) 非空。显然对 \(\sigma \in \mathcal{F}\)，有 \(\sigma \in \mathcal{G}^{\mathbb{Z}}(\mathfrak{p}')\) 且 \(\bar{\sigma} = u\)，这就在此情形完成了 (ii) 的证明。

\begin{enumerate}
\def\labelenumi{(\Alph{enumi})}
\setcounter{enumi}{2}
\tightlist
\item
  一般情形。G 的不变量域 \(K_{1}\) 是 K 的纯不可分扩张（《代数》第 V 章 § 10 第 9 节命题 14）；因此存在位于 p 之上的唯一素理想 \(p_{1}\)，即 \(A_{1} = A' \cap K_{1}\) 的素理想（引理 4）。若 \(p'\) 与 \(q'\) 是 \(A'\) 的位于 p 之上的两个素理想，则它们都位于 \(p_{1}\) 之上；由于 \(K'\) 是 \(K_{1}\) 的 Galois 扩张且 \(A' \cap K_{1}\) 整闭（§ 1 第 2 节命题 7 与命题 8 的推论），由 (B) 存在 \(\sigma \in G\) 使 \(\sigma \cdot p' = q'\)；由此得 (i)。另一方面，显然 \(k_{1}\)（即 \(A_{1}/p_{1}\) 的分式域）是 k 的纯不可分扩张（因 A 整闭）；由于 \(k'\) 由 (B) 是 \(k_{1}\) 的拟 Galois 扩张，故 \(k'\) 是 k 的拟 Galois 扩张，因为 \(k'\) 到 \(k'\) 的一个代数闭扩张的每个 k-同构都是 \(k_{1}\)-同构。最后这个注记结合 (B) 还表明，\(k'\) 的每个 k-自同构都具有形状 \(\bar{\sigma}\)，其中 \(\sigma \in \mathcal{G}^{\mathbb{Z}}(\mathfrak{p}')\)，这就完成了 (ii) 的证明。
\end{enumerate}

注

\begin{enumerate}
\def\labelenumi{(\arabic{enumi})}
\setcounter{enumi}{1}
\tightlist
\item
  设 \(\mathbf{K}'\) 是 \(\mathbf{K}\) 的 Galois 扩张，并保留命题 6 的证明中的记号；对每个 \(\alpha\)，设 \(\mathcal{G}_{\alpha}^{\mathbb{Z}}\)（相应地 \(\mathcal{G}_{\alpha}^{\mathrm{T}}\)）为 \(\mathcal{G}\) 的由诸 \(\sigma\)（其在 \(\mathrm{A}' \bigcap \mathrm{K}_{\alpha}\) 上的限制属于 \(\mathcal{G}^{\mathbb{Z}}(\mathfrak{p}' \bigcap \mathrm{K}_{\alpha})\)（相应地 \(\mathcal{G}^{\mathrm{T}}(\mathfrak{p}' \bigcap \mathrm{K}_{\alpha})\)）。命题 6 的证明表明这些子群在 \(\mathcal{G}\) 中闭，且
\end{enumerate}

\[
\mathcal {G} ^ {\mathrm{Z}} \left(\mathfrak {p} ^ {\prime}\right) = \bigcap_ {\alpha} \mathcal {G} _ {\alpha} ^ {\mathrm{Z}} \quad \text { and } \quad \mathcal {G} ^ {\mathrm{T}} \left(\mathfrak {p} ^ {\prime}\right) = \bigcap_ {\alpha} \mathcal {G} _ {\alpha} ^ {\mathrm{T}}.
\]

而且，在 \(A^{\prime} \cap K_{\alpha}\) 上，G（相应地 \(G_{\alpha}^{T}\)）的元素的限制所成的集就是整个群 \(\mathcal{G}^{\mathtt{Z}}(\mathfrak{p}^{\prime} \cap \mathtt{K}_{\alpha})\)（相应地 \(\mathcal{G}^{\mathtt{T}}(\mathfrak{p}^{\prime} \cap \mathtt{K}_{\alpha})\)），因为 \(K_{\alpha}\) 的每个 K-自同构都可延拓为 G 的一个元素。

在同一假设下，环 \(\mathbf{A}^{\mathbf{Z}}(\mathfrak{p}')\)（相应地 \(\mathbf{A}^{\mathbf{T}}(\mathfrak{p}')\)）是定向族 \(\mathbf{A}^{\mathbf{Z}}(\mathfrak{p}' \cap \mathbf{K}_{\alpha})\)（相应地 \(\mathbf{A}^{\mathbf{T}}(\mathfrak{p}' \cap \mathbf{K}_{\alpha})\)）的并：事实上，每个 \(x \in \mathbf{A}^{\mathbf{Z}}(\mathfrak{p}')\)（相应地每个 \(x \in \mathbf{A}^{\mathbf{T}}(\mathfrak{p}')\)）属于某个 \(\mathbf{K}_{\alpha}\)，且由前所述存在 \(\beta\) 使 \(\mathbf{K}_{\alpha} \subset \mathbf{K}_{\beta}\)，并且到 \(\mathbf{A}' \cap \mathbf{K}_{\alpha}\) 上，\(\mathcal{G}^{\mathbf{Z}}(\mathfrak{p}')\)（相应地 \(\mathcal{G}^{\mathbf{T}}(\mathfrak{p}')\)）的元素的限制与到 \(\mathbf{A}' \cap \mathbf{K}_{\alpha}\) 上 \(\mathcal{G}^{\mathbf{Z}}(\mathfrak{p}' \cap \mathbf{K}_{\beta})\)（相应地 \(\mathcal{G}^{\mathbf{T}}(\mathfrak{p}' \cap \mathbf{K}_{\beta})\)）的元素的限制相同，群 \(\mathcal{G}^{\mathbf{Z}}(\mathfrak{p}' \cap \mathbf{K}_{\alpha})\) 与 \(\mathcal{G}^{\mathbf{T}}(\mathfrak{p}' \cap \mathbf{K}_{\beta})\) 是有限的；因此 \(x\) 属于 \(\mathbf{A}^{\mathbf{Z}}(\mathfrak{p}' \cap \mathbf{K}_{\beta})\)（相应地 \(\mathbf{A}^{\mathbf{T}}(\mathfrak{p}' \cap \mathbf{K}_{\beta})\)）。

推论 1. 在命题 6 的假设下，设 \(f\) 是从 \(\mathbf{A}\) 到一个域 \(\mathbf{L}\) 的同态，\(g_1\)、\(g_2\) 是从 \(\mathbf{A}'\) 到 \(\mathbf{L}\) 且延拓 \(f\) 的两个同态。那么存在一个 \(\mathbf{K}\)-自同构 \(\sigma\)，即 \(\mathbf{K}'\) 的自同构，使 \(g_1 = g_2 \circ \sigma\)。

从命题 6 出发给出本推论的证明，与从定理 2 出发证明该定理的推论时的做法相同。

推论 2. 设 A 为整闭整环，K 为其分式域，K' 为 K 的有限代数扩张，A' 为 K' 中包含 A 且在 A 上整的子环。

\begin{enumerate}
\def\labelenumi{(\roman{enumi})}
\item
  对 \(\mathfrak{p}\)（\(\mathbf{A}\) 的每个素理想），\(\mathbf{A}'\) 的位于 \(\mathfrak{p}\) 之上的素理想所成的集是有限的。\\
\item
  若 \(\mathfrak{p}'\) 是 \(\mathbf{A}'\) 的位于 \(\mathfrak{p}\) 之上的一个素理想，则 \(\mathbf{A}' / \mathfrak{p}'\) 的每个元素都具有次数 \(\leqslant [\mathbf{K}' : \mathbf{K}]\)，这里是对 \(\mathbf{A} / \mathfrak{p}\) 的分式域而言。
\item
  设 \(\mathbf{K}''\) 为 \(\mathbf{K}\) 的拟 Galois 扩张，它是在 \(\mathbf{K}'\) 的一个代数闭包中由 \(\mathbf{K}'\) 生成的；设 \(\mathbf{A}''\) 为 \(\mathbf{A}\) 在 \(\mathbf{K}''\) 中的整闭包。域 \(\mathbf{K}''\) 是 \(\mathbf{K}\) 的有限扩张（《代数》第 V 章 § 6 第 3 节命题 9 的推论 1），因而其 K-自同构群有限；由此可知 \(\mathbf{A}''\) 的位于 \(p\) 之上的素理想所成的集是有限的（命题 6 (i)）。另一方面，由于 \(\mathbf{A}''\) 在 \(\mathbf{A}\) 上整，映射 \(p'' \mapsto p'' \cap \mathbf{A}'\)（从 \(\mathbf{A}''\) 的位于 \(p\) 之上的素理想之集到 \(\mathbf{A}'\) 的位于 \(x\) 之上的素理想之集）是满射（第 1 节定理 1）。
\item
  任意元素 \(x' \in A'\) 的极小多项式（在 K 上的）的系数都属于 A（§ 1 第 3 节命题 10 的推论）；对这个多项式的系数施加典范同态 \(\pi': A' \to A'/p'\)，便得到一个系数在 A/p 中、次数 \(\leqslant [K': K]\) 的整依赖方程，它为模 \(p'\) 之下的 \(x\) 的剩余类所满足；由此即得结论。
\end{enumerate}

推论 3. 在推论 2 的假设与记号下，若 A 是半局部环，则 A' 亦然。

对每个极大理想 \(m'\)（属于 \(A'\)），\(m' \cap A\) 都是 A 的一个极大理想（第 1 节命题 1）；由于假设 A 的极大理想所成的集有限，本推论随即由推论 2 得出。

推论 4. 设 A 为整闭整环，K 为其分式域，K' 为 K 的 Galois 扩张，A' 为 A 在 K' 中的整闭包，p' 为 A' 的一个素理想，p = A ∩ p'，且 k 与 k' 分别为 A/p 与 A'/p' 的分式域。那么：

\begin{enumerate}
\def\labelenumi{(\roman{enumi})}
\item
  \(\mathbf{A}^{\mathbf{z}} / (\mathfrak{p}'\cap \mathbf{A}^{\mathbf{z}})\) 的分式域等于 \(k\)，而 \(\mathbf{A}^{\mathbf{z}}\) 相对于 \(\mathfrak{p}'\cap \mathbf{A}^{\mathbf{z}}\) 的局部环的极大理想由 \(\mathfrak{p}\) 生成。
\item
  分式域 \(k^{\mathbf{T}}\)（即 \(\mathbf{A}^{\mathbf{T}} / (\mathfrak{p}' \cap \mathbf{A}^{\mathbf{T}})\) 的分式域）是 \(k\) 的包含在 \(k'\) 中的最大可分扩张。
\end{enumerate}

环 A 是 K' 在 K 上的 Galois 群在 A' 中的不变量环；若 K' 在 K 上是有限次的，则本推论由第 2 节的命题 4 与命题 5 得出。现在考虑一般情形，此时 K' 是由 K 的有限 Galois 扩张组成的右定向集 (K \(_{\alpha}\) ) 的并。那么：

\begin{enumerate}
\def\labelenumi{(\roman{enumi})}
\tightlist
\item
  若 \(x, y\) 是 \(\mathbf{A}^z\) 的两个元素，其中 \(y \notin \mathfrak{p}'\)，则存在指标 \(\alpha\) 使 \(x\) 与 \(y\) 属于 \(\mathbf{A}^z(\mathfrak{p}' \cap \mathbf{K}_\alpha)\)（注 2）；由第 2 节的命题 4，存在 \(x_0, y_0\)（在 \(\mathbf{A}\) 中），满足 \(y_0 \notin \mathfrak{p}'\) 且 \(xy_0 - x_0y \in \mathfrak{p}'\)，这就证明了 (i) 的第一个断言；若进一步 \(x \in \mathfrak{p}'\)，则可假定 \(y_0\) 满足
\end{enumerate}

\[
x y _ {0} \in \mathfrak {p A} ^ {\mathrm{z}} (\mathfrak {p} ^ {\prime} \cap \mathrm{K} _ {\alpha}) \subset \mathfrak {p A} ^ {\mathrm{z}} (\mathfrak {p} ^ {\prime}),
\]

这就证明了 (i) 的第二个断言。

\begin{enumerate}
\def\labelenumi{(\roman{enumi})}
\setcounter{enumi}{1}
\tightlist
\item
  现在设 \(x \in \mathbf{A}^{\mathrm{T}}\)；存在 \(\alpha\) 使 \(x \in \mathbf{A}^{\mathrm{T}}(\mathfrak{p}' \cap \mathbf{K}_{\alpha})\)（注 2），且第 2 节的命题 5 表明，剩余类 \(\bar{x}\)（即 \(x \bmod. (\mathfrak{p}' \cap \mathbf{K}_{\alpha} \cap \mathbf{A}^{\mathrm{T}})\) 的类）在 \(k\) 上是代数的且可分的；从而模 \((\mathfrak{p}' \cap \mathbf{A}^{\mathrm{T}})\) 之下 \(x\) 的剩余类更有理由在 \(k\) 上可分；为完成本推论的证明，只需证明 \(k'\) 是 \(k^{\mathrm{T}}\) 的纯不可分扩张。而 \(k'\) 是由分式域 \(k_{\alpha}\)（即环 \((\mathbf{A}' \cap \mathbf{K}_{\alpha}) / (\mathfrak{p}' \cap \mathbf{K}_{\alpha})\) 的分式域）组成的右定向族的并。于是由命题 5 可得：若 \(k'\) 的一个元素属于 \(k_{\alpha}\)，则它在如下商环的分式域上是纯不可分的
\end{enumerate}

\[
\mathbf {A} ^ {\mathrm{T}} \left(\mathfrak {p} ^ {\prime} \cap \mathbf {K} _ {\alpha}\right) / \left(\mathfrak {p} ^ {\prime} \cap \mathbf {A} ^ {\mathrm{T}} \left(\mathfrak {p} ^ {\prime} \cap \mathbf {K} _ {\alpha}\right)\right)
\]

从而更有理由在 \(k^{\mathbf{T}}\) 上是纯不可分的（依据注 2）。

定义 5. 在命题 6 的假设与记号下，不变量域 \(\mathbf{K}^{\mathrm{Z}}(\mathfrak{p}')\)（相应地 \(\mathbf{K}^{\mathrm{T}}(\mathfrak{p}')\)，即群 \(\mathcal{G}^{\mathrm{Z}}(\mathfrak{p}')\)（相应地 \(\mathcal{G}^{\mathrm{T}}(\mathfrak{p}')\)）在域 \(\mathbf{K}'\) 中的不变量域）称为 \(\mathfrak{p}'\) 关于 \(\mathbf{K}\) 的分解域（相应地惯性域）。

我们也用 \(K^{z}\)（相应地 \(K^{T}\)）代替 \(\mathbf{K}^{\mathbf{Z}}(\mathfrak{p}')\)（相应地 \(\mathbf{K}^{\mathrm{T}}(\mathfrak{p}')\)）来记。由 §1 第 9 节命题 23 可知，\(K^{z}\)（相应地 \(K^{T}\)）是环 \(A^{z}\)（相应地 \(A^{T}\)）的分式域；\(A^{z}\)（相应地 \(A^{T}\)）是 A 在 \(K^{z}\)（相应地 \(K^{T}\)）中的整闭包。

\textbf{注}\par

\begin{enumerate}
\def\labelenumi{(\arabic{enumi})}
\setcounter{enumi}{2}
\tightlist
\item
  在命题 6 的推论 4 的条件下，并设 \([\mathbf{K}':\mathbf{K}]\) 有限，位于 \(p\) 之上的不同素理想的个数是 \([\mathbf{K}^{\mathbf{Z}}:\mathbf{K}]\)，由 Galois 理论，此次数等于指数 \((\mathcal{G}:\mathcal{G}^{\mathbf{Z}})\)；此外，由 Galois 理论可得
\end{enumerate}

\[
[ \mathbf {K} ^ {\mathrm{T}}: \mathbf {K} ^ {\mathrm{Z}} ] = (\mathcal {G} ^ {\mathrm{Z}}: \mathcal {G} ^ {\mathrm{T}}) = [ k ^ {\mathrm{T}}: k ].\tag{6}
\]

\begin{enumerate}
\def\labelenumi{(\arabic{enumi})}
\setcounter{enumi}{3}
\tightlist
\item
  设 A 为整闭整环，K 为其分式域，\(K'\) 为 K 的有限代数扩张，\(A'\) 为 A 在 \(K'\) 中的整闭包。那么，对 A 的每个素理想 p，\(A'\) 的位于 p 之上的素理想的个数至多为 \([K':K]_{s}\)（即 \(K'\) 在 K 上的次数的可分部分）。我们可以先把注意力限制到 \(K'\) 是 K 的可分扩张的情形，因为一般地 \(K'\) 是 K 的最大可分扩张 \(K_{0}\)（即包含在 \(K'\) 中的那个）的纯不可分扩张，按定义 \([K':K]_{s} = [K_{0}:K]\)，而且若 \(A_{0}\) 是 A 在 \(K_{0}\) 中的整闭包，则 \(A_{0}\) 与 \(A'\) 的素理想一一对应（引理 4）。因此设 \(K'\) 在 K 上可分，并设 N 为在 K 的一个代数闭包中由 \(K'\) 生成的 K 的 Galois 扩张，G 为其 Galois 群，B 为 A 在 N 中的整闭包，P 为 B 的位于 p 之上的一个素理想。设 H 为 N 在 \(K'\) 上的 Galois 群，\(G^{z}\) 为 P 的分解群；B 的位于 p 之上的素理想就是诸 s.P，其中 \(s \in G\)（第 2 节定理 2），而关系 s.P = s'.P 意味着 \(s' = sg\)，其中 \(g \in G^{z}\)。另一方面，要使 s.P ∩ \(K' = s'\).P ∩ \(K'\)，必须且只需 \(s'\).P = ts.P，其中 \(t \in H\)（第 2 节定理 2），最终得 \(s' = tsg\)，其中 \(t \in H\) 且 \(g \in G^{z}\)。于是 A' 的位于 p 之上的素理想的个数等于 G 按等价关系「存在 \(t \in H\) 与 \(g \in G^{z}\) 使 \(s' = tsg\)」（在 s 与 \(s'\) 之间）所分的类的个数；显然这个数至多等于指数 \((\mathcal{G}:\mathcal{H})\)，即 H 在 G 中的右陪集数，而由 Galois 理论 \((\mathcal{G}:\mathcal{H}) = [\mathrm{K}':\mathrm{K}]\)。
\end{enumerate}

命题 7. 设 A 为整闭整环，K 为其分式域，K' 为 K 的 Galois 扩张，G 为其 Galois 群，A' 为 A 在 K' 中的整闭包，p' 为 A' 的素理想，且 p = A ∩ p'。最后，设 L 为 K' 中包含 K 的子域，并令 p(L) = p' ∩ L。

\begin{enumerate}
\def\labelenumi{(\roman{enumi})}
\item
  \(\mathfrak{p}'\) 关于 \(\mathbf{L}\) 的分解域（相应地惯性域）是 \(\mathbf{L}(\mathbf{K}^{\mathbf{Z}})\)（相应地 \(\mathbf{L}(\mathbf{K}^{\mathbf{T}})\)）；若进一步 \(\mathbf{L}\) 是 \(\mathbf{K}\) 的 Galois 扩张，则 \(\mathfrak{p}(\mathbf{L})\) 关于 \(\mathbf{K}\) 的分解域是 \(\mathbf{L} \cap \mathbf{K}^{\mathbf{Z}}\)。
\item
  若 \(\mathbf{L}\) 包含于 \(\mathbf{K}^{\mathbf{Z}}\)，则 \(\mathbf{A} / \mathfrak{p}\) 与 \((\mathbf{A}' \cap \mathbf{L}) / \mathfrak{p}(\mathbf{L})\) 有相同的分式域，并且在局部环 \(\mathbf{A}' \cap \mathbf{L}\)（对应于素理想 \(\mathfrak{p}(\mathbf{L})\)）中，极大理想由 \(\mathfrak{p}\) 生成。反之，若这两个条件成立，且进一步 \(\bigcap_{n \geqslant 0} \mathfrak{p}^{n} \mathbf{A}_{\mathfrak{p}'}' = 0\)，则 \(\mathbf{L}\) 包含于 \(\mathbf{K}^{\mathbf{Z}}\)。
\item
  若 \(\mathcal{H}\) 是 \(\mathcal{G}\) 的使 \(\mathbf{L}\) 的元素保持不变的那个子群，则显然 \(\mathfrak{p}'\) 关于 \(\mathbf{L}\) 的分解群（相应地惯性群）是 \(\mathcal{G}^{\mathrm{Z}}\cap \mathcal{H}\)（相应地 \(\mathcal{G}^{\mathrm{T}}\cap \mathcal{H}\)），而当 \(\mathbf{K}'\) 是 \(\mathbf{K}\) 的有限 Galois 扩张时，第一个断言由 Galois 理论得出（《代数》第 V 章 § 10 第 6 节定理 3 的推论 1）；一般情形则由如下事实得出：在注 2 的记号下，\(\mathbf{A}^{\mathrm{Z}}\)（相应地 \(\mathbf{A}^{\mathrm{T}}\)）是诸 \(\mathbf{A}^{\mathrm{Z}}(\mathfrak{p}'\cap \mathbf{K}_{\alpha})\)（相应地 \(\mathbf{A}^{\mathrm{T}}(\mathfrak{p}'\cap \mathbf{K}_{\alpha})\)）的并：每个元素 \(x\in \mathbf{K}'\) 属于某个 \(\mathbf{K}_{\alpha}\)，而若它在 \(\mathcal{G}^{\mathrm{Z}}(\mathfrak{p}')\cap \mathcal{H}\)（相应地 \(\mathcal{G}^{\mathrm{T}}(\mathfrak{p}')\cap \mathcal{H}\)）下不变，则它也在 \(\mathcal{G}^{\mathrm{Z}}(\mathfrak{p}'\cap \mathbf{K}_{\beta})\cap \mathcal{H}\)（相应地 \(\mathcal{G}^{\mathrm{T}}(\mathfrak{p}'\cap \mathbf{K}_{\beta})\cap \mathcal{H}\)）下不变，其中 \(\beta\) 是某个适当的指标；于是由论证的开头部分，它属于
\end{enumerate}

\[
\mathbf {L} \left(\mathrm{K} ^ {\mathrm{Z}} \left(\mathfrak {p} ^ {\prime} \cap \mathrm{K} _ {\alpha}\right)\right) \subset \mathbf {L} \left(\mathrm{K} ^ {\mathrm{Z}}\right) (\text { resp. } \mathbf {L} \left(\mathrm{K} ^ {\mathrm{T}} \left(\mathfrak {p} ^ {\prime} \cap \mathrm{K} _ {\beta}\right)\right) \subset \mathbf {L} \left(\mathrm{K} ^ {\mathrm{T}}\right)).
\]

现在设 L 是 K 的 Galois 扩张；于是每个 \(\sigma\in G^{z}\) 在 L 上的限制都保持 \(\mathfrak{p}(L)=\mathfrak{p}^{\prime}\cap L\) 不变，从而属于 \(\mathfrak{p}(L)\) 关于 K 的分解群。反之，设 \(\tau\) 是属于该群的 L 的一个自同构，并设 \(\sigma\) 是 \(\tau\) 到 \(K^{\prime}\) 的一个 K-自同构上的延拓；记 \(q^{\prime}=\sigma.p^{\prime}\)。由于 \(p^{\prime}\) 与 \(q^{\prime}\) 都位于 \(\mathfrak{p}(L)\) 之上，故存在自同构 \(\rho\in H\) 使 \(q^{\prime}=\rho.p^{\prime}\)，从而 \(\rho^{-1}\sigma\in G^{z}\)，且 \(\tau\) 是 \(\rho^{-1}\sigma\) 在 L 上的限制；换句话说，\(\mathfrak{p}(L)\) 关于 K 的分解群恰为诸自同构 \(\sigma\in G^{z}\) 在 L 上的限制所成的群，这就证明了第二个断言。

\begin{enumerate}
\def\labelenumi{(\roman{enumi})}
\setcounter{enumi}{1}
\tightlist
\item
  说 \(L \subset K^{z}\) 即意味着 \(H \supset G^{z}\)，因此当 \([K':K]\) 有限时，(ii) 的断言是第 2 节命题 4 (ii) 与 (iii) 的特殊情形。一般情形下的论证与命题 6 的证明中的相同。
\end{enumerate}

